\chapter{Direzioni future}
\label{ch:direzioni-future}

Questo capitolo delinea le direzioni di sviluppo più promettenti, scientifiche e applicative: l'applicativo fornisce l'infrastruttura di dati e validazione su cui le domande di ricerca possono essere testate empiricamente.

\section{Direzioni per la ricerca}

\subsection{Stima empirica del parametro di decay}

Nel prototipo l'emivita del relevance score è configurabile (default: 7 giorni) ma non stimata dai dati.
Una prima domanda a basso costo sperimentale è la stima ottimale di $h$, selezionando il valore che massimizza la correlazione ritardata tra KPI di sentiment e rendimenti successivi, per titolo o settore.
L'ipotesi verificabile è che l'emivita informativa non sia costante, ma vari con capitalizzazione, copertura mediatica e regime di mercato---un risultato autonomo rispetto al decay fisso dei provider commerciali.

\subsection{Diffusione su grafo: modello discreto e garanzie di correttezza}
\label{sec:soluzione-discreta}

La formulazione a derivate parziali della diffusione semantica (Capitolo~\ref{ch:descrizione-progetto}) diventa operativa nella sua approssimazione discreta su grafo, di cui il prototipo possiede già l'ingrediente chiave: la matrice di similarità tra profili aziendali.
Rispetto alla letteratura su sentiment spillover e graph neural network (Capitolo~\ref{ch:stato-dell-arte}), il vantaggio distintivo è l'interpretabilità---$D$ e $\lambda$ sono parametri fisicamente leggibili, a differenza dei pesi appresi da una GNN---e il confronto sistematico tra diffusione interpretabile e baseline GNN costituirebbe il contributo scientifico centrale del progetto.
Questa sezione formula il modello discreto completo, incluso il trasporto del sentiment la cui necessità è discussa nel Capitolo~\ref{ch:descrizione-progetto}, e ne dimostra le proprietà di correttezza.

\paragraph{Modello.}
Sia $G$ un grafo non orientato su $n$ nodi (i ticker), con matrice dei pesi $W \in \mathbb{R}^{n \times n}$ simmetrica, a elementi non negativi e diagonale nulla (similarità coseno tra embedding, sogliata per sparsificare).
Siano $d_i = \sum_j w_{ij}$ i gradi, $D_g = \operatorname{diag}(d_1, \dots, d_n)$ e $L = D_g - W$ il Laplaciano combinatorio.
Lo stato è la coppia $(R(t), m(t))$, dove $R_i$ è la rilevanza attiva sul nodo $i$ e $m_i$ la sua \emph{massa di sentiment}, inizializzata come $m_i(0) = S_i(0)\, R_i(0)$; entrambe evolvono con lo stesso operatore:
\[
\frac{dR}{dt} = -(D\,L + \lambda I)\, R,
\qquad
\frac{dm}{dt} = -(D\,L + \lambda I)\, m,
\qquad
S_i(t) := \frac{m_i(t)}{R_i(t)} \quad (R_i(t) > 0),
\]
con $D \geq 0$ coefficiente di diffusione e $\lambda > 0$ tasso di decadimento: la rilevanza in movimento trasporta il proprio tono.

\begin{proposizione}[Esistenza, unicità e forma chiusa]
\label{prop:esistenza}
Il sistema ammette un'unica soluzione,
\[
R(t) = e^{-\lambda t}\, e^{-D t L}\, R(0),
\qquad
m(t) = e^{-\lambda t}\, e^{-D t L}\, m(0);
\]
se $L = U \Lambda U^{\top}$ è la decomposizione spettrale, $e^{-DtL} = U e^{-Dt\Lambda} U^{\top}$.
\end{proposizione}
\begin{proof}
Sistema lineare a coefficienti costanti: la soluzione è l'esponenziale di matrice $e^{-t(DL + \lambda I)}$, che fattorizza perché $\lambda I$ commuta con $L$.
La decomposizione spettrale esiste per il teorema spettrale ($L$ simmetrica semidefinita positiva); il costo $O(n^3)$, una tantum, è trattabile per $n \sim 2 \cdot 10^3$.
\end{proof}

\begin{proposizione}[Positività]
\label{prop:positivita}
Se $R(0) \geq 0$ componente per componente, allora $R(t) \geq 0$ per ogni $t \geq 0$.
\end{proposizione}
\begin{proof}
Sia $c \geq \max_i d_i$: la matrice $cI - L = W + (cI - D_g)$ ha elementi non negativi, e
$e^{-DtL} = e^{-Dtc}\, e^{Dt(cI - L)}$,
dove l'esponenziale di una matrice a elementi non negativi è a elementi non negativi (serie di potenze a termini non negativi).
Dunque $e^{-DtL} \geq 0$ e $R(t) \geq 0$.
\end{proof}

\begin{proposizione}[Bilancio di massa e compatibilità retroattiva]
\label{prop:massa}
$\sum_i R_i(t) = e^{-\lambda t} \sum_i R_i(0)$: la diffusione redistribuisce la rilevanza senza crearne né distruggerne; solo $\lambda$ la dissipa.
Per $D = 0$ il modello si riduce a $R_i(t) = R_i(0) \cdot 0{,}5^{\,t/h}$ con $h = \ln 2 / \lambda$: il sistema in produzione è il caso particolare $D = 0$.
\end{proposizione}
\begin{proof}
$L \mathbf{1} = 0$ per costruzione e, per simmetria, $\mathbf{1}^{\top} L = 0$; quindi $\frac{d}{dt}(\mathbf{1}^{\top} R) = -\lambda\, \mathbf{1}^{\top} R$.
Per $D = 0$ le equazioni si disaccoppiano in $R_i' = -\lambda R_i$, che riscritta in emivita coincide con il peso $w = 0{,}5^{\Delta t / h}$ del Capitolo~\ref{ch:stato-implementazione}.
\end{proof}

La compatibilità retroattiva è anche il test di validazione naturale: azzerando $W$, la nuova pipeline deve riprodurre esattamente i KPI attuali.

\begin{proposizione}[Principio del massimo per il sentiment]
\label{prop:massimo}
Per ogni nodo $i$ e $t \geq 0$ con $R_i(t) > 0$:
$\min_j S_j(0) \leq S_i(t) \leq \max_j S_j(0)$.
\end{proposizione}
\begin{proof}
Sia $P(t) = e^{-t(DL + \lambda I)} \geq 0$ (Proposizione~\ref{prop:positivita}).
Allora
\[
S_i(t) = \frac{(P(t)\, m(0))_i}{(P(t)\, R(0))_i}
       = \frac{\sum_j P_{ij}(t)\, R_j(0)\, S_j(0)}{\sum_j P_{ij}(t)\, R_j(0)}
       = \sum_j \alpha_j S_j(0),
\]
con $\alpha_j \geq 0$ e $\sum_j \alpha_j = 1$: combinazione convessa dei sentiment iniziali.
\end{proof}

È la garanzia applicativa decisiva: la diffusione \emph{non può fabbricare} un sentiment più estremo di quello osservato nelle notizie---nessuna amplificazione spuria in un sistema che genera segnali operativi.

\begin{proposizione}[Schema numerico: stabilità e positività]
\label{prop:schema}
Lo schema di Eulero esplicito $R_{k+1} = (I - \Delta t\,(DL + \lambda I)) R_k$ è consistente di ordine~1 e, se
$\Delta t \leq 1/(\lambda + D\, d_{\max})$ con $d_{\max} = \max_i d_i$,
preserva positività e stabilità.
Lo schema implicito $R_{k+1} = (I + \Delta t\,(DL + \lambda I))^{-1} R_k$ è incondizionatamente stabile e positivo.
\end{proposizione}
\begin{proof}
\emph{Esplicito.} La matrice di iterazione $M = I - \Delta t (DL + \lambda I)$ ha fuori diagonale $\Delta t D\, w_{ij} \geq 0$ e diagonale $1 - \Delta t (D d_i + \lambda) \geq 0$ sotto la condizione: $M \geq 0$.
Gli autovalori sono $1 - \Delta t(\lambda + D\mu)$ con $\mu \in [0, \lambda_{\max}(L)]$ e $\lambda_{\max}(L) \leq 2 d_{\max}$; sotto la condizione, $\Delta t (\lambda + D \lambda_{\max}) \leq (\lambda + 2 D d_{\max})/(\lambda + D d_{\max}) \leq 2$, quindi lo spettro resta in $[-1, 1]$; la convergenza $O(\Delta t)$ è classica.
\emph{Implicito.} $I + \Delta t (DL + \lambda I)$ ha diagonale strettamente dominante e fuori diagonale non positiva: è una M-matrice non singolare con inversa $\geq 0$, e il raggio spettrale dell'iterazione è $\leq (1 + \lambda \Delta t)^{-1} < 1$.
\end{proof}

\begin{osservazione}[Calibrazione]
Con passo giornaliero e l'emivita di default ($h = 7$, $\lambda \approx 0{,}099$), la condizione diventa $D\, d_{\max} \lesssim 0{,}9$: un vincolo interpretabile tra intensità della diffusione e densità del grafo, controllabile con la soglia di similarità.
Lo schema implicito (una soluzione di sistema sparso al giorno) rimuove ogni vincolo sul passo.
\end{osservazione}

Il popolamento di \texttt{ticker\_cross\_diffused\_sentiment} si riduce quindi a un passo di Eulero implicito al giorno sulle coppie $(R, m)$: nessun ingrediente mancante.

\subsection{Benchmark multi-scorer del sentiment}

La delega dello scoring ad Alpha Vantage può diventare un'opportunità: il database contiene decine di migliaia di coppie (articolo, ticker) già etichettate.
Si propone un confronto sistematico tra scorer sullo stesso corpus---FinBERT zero-shot, un LLM open-weight via prompting, i punteggi del provider---su due domande: la concordanza inter-rater e la validità predittiva rispetto ai rendimenti.
La conversione \textit{logit-to-score} di FinDPO \parencite{iacovides2025findpo} offre punteggi continui e calibrati, confrontabili con la scala $[-1, +1]$ in uso.

\subsection{Attenzione mediatica come segnale complementare}

Il volume di notizie per ticker e topic, già raccolto dalla coverage ingestion, permette di costruire un indice di attenzione mediatica nello spirito dell'Hype Index \parencite{cao2025hypeindexnlpdrivenmeasure} senza costi aggiuntivi.
L'ipotesi è che i picchi di copertura anticipino volatilità indipendentemente dalla polarità, e che l'interazione sentiment~$\times$~attenzione superi i due segnali separati.

\subsection{Interazione tra sentiment e segnali tecnici}

Con gli indicatori tecnici (Sezione~\ref{sec:analisi-tecnica}) il prototipo espone per ogni titolo un segnale informativo e uno di prezzo; la domanda naturale è se il sentiment abbia \emph{valore predittivo incrementale} rispetto ai segnali tecnici, e se le divergenze anticipino inversioni.
\textcite{neely2014forecasting} offrono il precedente diretto: gli indicatori tecnici prevedono il premio al rischio con potere comparabile e \emph{complementare} alle variabili macroeconomiche.
Il disegno proposto replica quello schema sostituendo alle variabili macro il sentiment: regressioni dei rendimenti su sentiment, segnale tecnico e interazione \parencite{brock1992simple}.
Il report automatico del Capitolo~\ref{ch:stato-implementazione} già classifica ogni titolo come allineato o divergente: le etichette costituiscono le celle del disegno sperimentale, a costo quasi nullo.

\subsection{Notizie programmate e non programmate}

\textcite{boudoukh2019information} mostrano che l'impatto delle notizie firm-specific differisce tra eventi \emph{programmati} (earnings, guidance) e \emph{non programmati} (M\&A, azioni legali).
La pipeline già classifica gli articoli per topic (inclusi \texttt{earnings} e \texttt{mergers\_and\_acquisitions}): separare il KPI nei due regimi permetterebbe di verificare se l'emivita informativa $h$ differisca tra notizie attese e sorprese, raffinando la stima empirica del decay proposta sopra---un'analisi resa possibile dallo storico pluriennale del cold store.

\subsection{Foundation model per serie temporali}

La proposta originaria individua nell'LSTM il modello predittivo di riferimento; dal 2024 i modelli fondazionali per serie temporali---Chronos \parencite{ansari2024chronos}, TimesFM \parencite{das2024timesfm}---raggiungono in zero-shot prestazioni comparabili ai modelli supervisionati per singolo dataset.
La revisione a basso costo del piano: usare un foundation model come baseline zero-shot e verificare se il condizionamento sul sentiment (covariata o fine-tuning leggero) apporti valore incrementale.
Il confronto ``LSTM da zero vs.\ foundation model condizionato sul sentiment'' è a sua volta una domanda di ricerca attuale.

\subsection{Protocollo di backtesting metodologicamente rigoroso}

La letteratura critica \parencite{li2025finsaber} mostra che i vantaggi delle strategie LLM svaniscono su orizzonti lunghi e universi ampi, e che i backtest sono esposti a survivorship bias, look-ahead bias e data-snooping.
Il backtesting sarà quindi vincolato a un protocollo esplicito: solo informazione disponibile alla data di segnale, benchmark passivi, più regimi di mercato, costi di transazione dichiarati.
Per gli scorer LLM va considerato anche il look-ahead bias ``parametrico'': la valutazione pulita richiede notizie successive al cutoff di addestramento del modello.

\section{Direzioni per l'applicativo}

Le priorità implementative, per rapporto valore/sforzo:

\begin{enumerate}
  \item \textbf{Scoring interno con FinBERT}: applicazione batch agli articoli in database, con punteggi persistiti accanto a quelli del provider---abilita il benchmark multi-scorer;
  \item \textbf{Popolamento del cross-diffused sentiment} con il modello della Sezione~\ref{sec:soluzione-discreta} e la matrice di correlazione esistente;
  \item \textbf{Analisi statistica sentiment--prezzo}: correlazioni di Pearson e Spearman con ritardi configurabili ($0$--$5$ giorni), al posto del confronto visivo;
  \item \textbf{Backtester elementare}: strategia long/flat sui dati persistiti, con rendimento cumulato, hit-rate e Sharpe semplificato, nel rispetto del protocollo anti-bias;
  \item \textbf{Seconda fonte di dati} a basso attrito (RSS o scraping mirato); gli articoli senza scoring del provider sono il banco di prova del modulo FinBERT;
  \item \textbf{LLM-as-annotator} per un dataset entity-level nello spirito di FinEntity \parencite{tang-etal-2023-finentity}, con revisione manuale a campione, per abilitare il fine-tuning a costi contenuti.
\end{enumerate}

Il vincolo pratico più stringente---uno storico ampio su cui addestrare e validare---è stato rimosso dal cold store (Sezione~\ref{sec:cold-storage}), che accumula senza limiti di ritenzione il flusso di articoli etichettati dal 2018: fine-tuning e benchmark possono essere pianificati su centinaia di migliaia di esempi.

\section{Sintesi}

Il filo conduttore è la trasformazione del prototipo da strumento di monitoraggio a piattaforma sperimentale: ogni componente applicativa corrisponde a una domanda di ricerca verificabile con i dati già raccolti.
L'elemento di maggiore originalità resta la modellazione della rilevanza informativa come campo dinamico che decade e si diffonde nello spazio semantico---ora dotata di formulazione discreta, soluzione in forma chiusa e garanzie di correttezza dimostrate.
